\begin{semiwebaddition}
\section{Pauli semiwebs in double checking}
\label{app:semiwebs}

\subsection{Definition and relation to Pauli propagation}
A Pauli semiweb assigns $I,X,Y$ or $Z$ to every wire of a ZX diagram.
At each spider, the tensor product $w_v$ of the Pauli operators on its legs acts on its tensor as
\begin{equation}
 w_v|v(\theta_v)\rangle=\lambda_v|v(\theta_v+\delta_v)\rangle.
 \label{eq:semiweb-local}
\end{equation}
Here $|v(\theta_v)\rangle$ is the spider tensor written as a state,
$\lambda_v$ is a nonzero scalar, and a phase shift $\delta_v\ne0$
modulo $2\pi$ is a \emph{defect}. The Pauli operators on each Hadamard tensor's legs must leave it unchanged
up to a scalar. With no defects, the semiweb is a \emph{Pauli web}~\cite{zxflow}.
A non-Clifford diagram can have both Pauli webs and semiwebs with defects.
The accompanying \texttt{pauli\_semiwebs.pdf} gives a basis of the full
semiweb space for each $d=3,5$ double-checking diagram, with editable TikZ sources.

The relation to Pauli propagation is already apparent with only a single $T$ gate:
\begin{equation}
 T^\dagger X=e^{-i\pi/4}XT,
 \qquad T^\dagger XT=\frac{X-Y}{\sqrt2}.
 \label{eq:semiweb-T}
\end{equation}
The first identity keeps the Pauli label and changes the gate's phase;
the second expands the conjugated observable into signed Pauli terms.
Both retain the coherent phases.

Let $K^{(3)}$ represent the noiseless $d=3$ double-checking stage with every $X$
measurement selected as $+1$, before the terminal code projection.
The local identities in Eq.~\eqref{eq:semiweb-local} give
\begin{equation}
 K^{(3)} P_{\rm in}=\lambda P_{\rm out}K^{(3)}_{\boldsymbol\delta},
 \label{eq:semiweb-global}
\end{equation}
where $K^{(3)}_{\boldsymbol\delta}$ has the starred phases changed and $\lambda$
includes the scalar phases from the local identities and contractions.
For Fig.~\ref{fig:semiweb-main}, this is
\begin{equation}
 K^{(3)}X_7=e^{5i\pi/4}X_0X_3X_9X_{10}X_{12}X_{13}K^{(3)}_{\boldsymbol\delta}.
 \label{eq:semiweb-X7}
\end{equation}
The diagram describes the selected branch. It shows changes to phases and measurement outcomes, including cancellations;
full post-selection must still be checked.

\subsection{Multiplication and cancellation}
Products of semiwebs on the same diagram are semiwebs~\cite{zxflow}.
To multiply two semiwebs, multiply their Pauli operators on each wire, including scalar phases such as
$XZ=-iY$. Defect shifts must be composed, rather than added.

\paragraph{A non-Clifford defect that cancels.}
Both $W_{X_7}$ and $W_{X_0}$ put $X$ on both legs of the
$\pi/4$ spider at the exit of qubit $0$. Their product gives $XX=I$,
as shown in Fig.~\ref{fig:semiweb-red-local}. Its phase changes as
\begin{equation}
 \frac{\pi}{4}\xrightarrow{X\otimes X}-\frac{\pi}{4}
 \xrightarrow{X\otimes X}\frac{\pi}{4}.
 \label{eq:semiweb-phase-cancellation}
\end{equation}
The second change is $+\pi/2$ after the first has acted. Other defects remain.

\begin{figure}[p]
\centering
{\color{black}\textbf{(a)} $W_{X_7}$: incoming $X_7$.}\par\vspace{1mm}
\semiwebinput[.91\linewidth]{figures/semiwebs/x7_cancellation.tikz}
\par\vspace{1mm}
{\color{black}\textbf{(b)} $W_{X_0}$: incoming $X_0$.}\par\vspace{1mm}
\semiwebinput[.91\linewidth]{figures/semiwebs/x0_cancellation.tikz}
\par\vspace{1mm}
{\color{black}\textbf{(c)} Product $W_{X_7}W_{X_0}$: incoming $X_0X_7$.}\par\vspace{1mm}
\semiwebinput[.91\linewidth]{figures/semiwebs/x0x7_cancellation.tikz}
\semiwebcaption{Cancellation of a shared non-Clifford defect on the complete
$d=3$ double-checking circuit. Red overlays show $X$ Pauli semiwebs; stars mark
defects. The dashed circle picks out the same $T$-gate at qubit $0$ in all three
panels. In (c), $XX=I$ on each leg and the shared defect cancels, while other defects
remain.}
\label{fig:semiweb-red-local}
\end{figure}


\paragraph{Two measurement defects that cancel.}
The $Z_0$ and $Z_3$ semiwebs each flip the central $X$ measurement on
qubit $7$ and the final ancilla $X$ measurement on qubit $6$.
Their product cancels both measurement defects,
Fig.~\ref{fig:semiweb-cancellation}:
\begin{equation}
 W_{Z_0}W_{Z_3}=W_{Z_0Z_3},\qquad K^{(3)}Z_0Z_3=Z_0Z_3K^{(3)}.
 \label{eq:semiweb-Z-product}
\end{equation}
This product is a defect-free Pauli web.

\begin{figure}[p]
\centering
{\color{black}\textbf{(a)} Incoming $Z_0$: two measurement defects, $\delta=\pi$.}\par\vspace{1mm}
\semiwebinput{figures/semiwebs/z0.tikz}
\par\vspace{1mm}
{\color{black}\textbf{(b)} Incoming $Z_3$: the same two measurement defects.}\par\vspace{1mm}
\semiwebinput{figures/semiwebs/z3.tikz}
\par\vspace{1mm}
{\color{black}\textbf{(c)} Product $W_{Z_0}W_{Z_3}$: the shared $Z$ operators and measurement defects cancel.}\par\vspace{1mm}
\semiwebinput{figures/semiwebs/z0z3.tikz}
\semiwebcaption{Cancellation of the two measurement defects.
Green overlays show $Z$ Pauli semiwebs; violet stars mark $\pi$ defects in
selected measurement outcomes.}
\label{fig:semiweb-cancellation}
\end{figure}


\FloatBarrier
\end{semiwebaddition}
